\documentclass[11pt,a4paper]{article}

\usepackage[utf8]{inputenc}
\usepackage[T1]{fontenc}
\usepackage{lmodern}
\usepackage{microtype}
\usepackage[margin=2.5cm]{geometry}
\usepackage{amsmath,amssymb,amsthm}
\usepackage{mathtools}
\usepackage{booktabs}
\usepackage{array}
\usepackage{enumitem}
\usepackage{hyperref}
\usepackage{xcolor}
\usepackage{tcolorbox}
\usepackage{fancyhdr}

\hypersetup{colorlinks=true, linkcolor=blue!50!black,
  citecolor=blue!50!black, urlcolor=blue!50!black}

\pagestyle{fancy}
\fancyhf{}
\fancyhead[L]{\small A conditional substrate localisation of Hilbert--P\'olya}
\fancyhead[R]{\small Pre-peer-review preprint}
\fancyfoot[C]{\thepage}

\newtheorem{theorem}{Theorem}
\newtheorem{proposition}[theorem]{Proposition}
\newtheorem{lemma}[theorem]{Lemma}
\newtheorem{corollary}[theorem]{Corollary}
\newtheorem{definition}[theorem]{Definition}
\newtheorem{target}[theorem]{Target}
\newtheorem{remark}[theorem]{Remark}
\newtheorem{conjecture}[theorem]{Conjecture}

\newcommand{\II}{\mathcal{I}}
\newcommand{\Vsix}{V_{600}}
\newcommand{\Cphi}{C_\varphi}
\newcommand{\NH}{N_H}
\newcommand{\Q}{\mathbb{Q}}
\newcommand{\Z}{\mathbb{Z}}
\newcommand{\R}{\mathbb{R}}
\newcommand{\C}{\mathbb{C}}
\newcommand{\N}{\mathbb{N}}
\newcommand{\Vmin}{V_{\min}}
\newcommand{\Lsub}{L_{\mathrm{sub}}}
\newcommand{\Wit}{W}
\newcommand{\Twit}{T_W}

\newtcolorbox{scopebox}[1][]{
  colback=gray!3, colframe=gray!50!black,
  fonttitle=\bfseries, title=Scope, #1}

\newtcolorbox{equivbox}[1][]{
  colback=red!2, colframe=red!50!black,
  fonttitle=\bfseries, title={The localisation, stated precisely}, #1}

\title{\textbf{A Conditional Substrate Localisation of the
Hilbert--P\'olya Problem}\\[0.4em]
\large A target specification on an infinite-dimensional lift of the
$\tau$-fixed block of $V_{600}$: if a generator with the stated
substrate constraints exists, its positive eigenvalues in
multiplicity-respecting bijection with $\zeta$'s upper-half-plane
non-trivial zeros, then RH follows; the converse is the named open
problem}

\author{Lee Smart\\
\small Institute of Vibrational Field Dynamics\\
\small \texttt{contact@vibrationalfielddynamics.org}}

\date{2026-06-12 \quad (pre-peer-review preprint, v2.0.0-rc1 --- reframed from v1 after adversarial review)}

\begin{document}
\maketitle

\begin{abstract}
\noindent
The Riemann Hypothesis (RH) is the conjecture that every non-
trivial zero of $\zeta(s)$ has $\mathrm{Re}(s) = 1/2$. The
$\Vsix$ closure programme has established, in eight findings
\cite{criticalLinePullback}, a substrate-side framework in which:
the critical line embeds into $\Vsix$ via the Hilbert modular
L-function $\Lsub = \zeta_K \cdot \zeta_K(s-1) \cdot C_2$;
the closure operator $\Cphi$ has a Galois decomposition $94 + 13 +
13$; and the witness invariant $\Wit = \Vmin$ is the universal
attractor of closure-flow \cite{observerAttractor}.

This paper states a \emph{conditional localisation} of the
Hilbert--P\'olya problem onto that framework.

\medskip
\noindent
\textbf{Proposition (one direction, classical).}
\itshape
If a self-adjoint operator $\Twit$ exists on an
infinite-dimensional lift $\mathcal{H}_+$ of the $\tau$-fixed
block of $\Vsix$, satisfying the substrate constraints of
Definition~\ref{def:spectral-generator}, whose positive eigenvalues are in multiplicity-respecting
bijection with the non-trivial zeros of $\zeta$ in the upper half
plane (via $\rho = \tfrac12 + i\,t_\rho$; the conjugate lower
half follows by symmetry), then RH holds.

\medskip
\noindent
\upshape
This direction is the classical Hilbert--P\'olya implication
(self-adjointness forces a real spectrum) and carries no novelty;
what is ours is the \emph{constraint list} --- ground state at
$\Vmin$, commutation with closure-flow and with the Hecke action
of the Hilbert modular lift --- which localises \emph{where} such
an operator would have to live. Three honesty notes, up front.
First, no lift $\mathcal{H}_+$ is constructed here: the
$\tau$-fixed block is $94$-dimensional, a finite space cannot
carry the infinite spectrum, and producing a lift with the stated
properties is the open problem this paper names. Second, the
converse --- that RH guarantees a generator satisfying the
substrate constraints --- is \emph{not} established here, and
without the constraints the converse is trivial (a diagonal
operator on $\ell^2$ does it), so the entire content of the
localisation lies in the constraints. Third, this is a
specification, not a theorem about $\zeta$: it sharpens the
question, it does not move the answer.

\bigskip

\noindent\textbf{Status.} Pre-peer-review open research preprint.
Not independently validated. No claim of RH settlement is made.

\end{abstract}

\bigskip

\tableofcontents

\newpage

% ============================================================
\section{Why RH is the final reframe}
\label{sec:intro}

The $\Vsix$ closure programme has produced a structural picture
in which the Riemann critical line has a precise substrate
interpretation. The Witness Invariant Theorem
\cite{observerAttractor} establishes:

\begin{itemize}[leftmargin=*]
  \item the substrate is the $600$-cell vertex set $\Vsix$;
  \item the closure equation
        $\Wit = \Phi \times G \times O$ subject to $C \to 0$
        has the unique solution $\Wit = \Vmin$;
  \item every closure-flow trajectory attracts to $\Vmin$;
  \item the critical line embeds into the $\tau$-fixed $94$
        block via $\Lsub$.
\end{itemize}

What remains open is the analytical step: are the ordinates
$\gamma_n$ along the critical line $s = 1/2 + it$
\emph{forced} to lie at the eigenvalues of some specific operator
on the substrate's $\tau$-fixed block? This is the substrate-side
version of the Hilbert--P\'olya conjecture.

\subsection{The role of RH in the puzzle}

The Witness Invariant Theorem identifies \emph{where} the witness
lives ($\Vmin$). It does not specify \emph{which excited
eigenvalues} the substrate's spectral generator has. RH is the
statement that constrains these eigenvalues to be exactly
$\{\gamma_n\}$.

In the diagram's vocabulary (\cite{observerAttractor}, Fig.~1):
\begin{itemize}[leftmargin=*]
  \item ``balance $1/2$ holds'' --- this is the witness axis
        $\mathrm{Re}(s) = 1/2$, fixed by the functional equation;
  \item ``$t$ unfolds'' --- the spectral parameter $t$ along the
        critical line is the parameter that gives the spectral
        excitation;
  \item the ordinates $\gamma_n$ are the $t$-values
        where $\zeta$ vanishes on the line (under RH, all of them).
\end{itemize}

RH says these $t$-values are real and lie at specific positions.
The substrate framework identifies $\Vmin$ as ``$t = 0$, balance''
and asks whether the excited eigenstates have $t$-values matching
the $\gamma_n$.

\subsection{What this paper does}

We articulate the precise sense in which a substrate-localised
Hilbert--P\'olya operator would witness RH. The main statement
(\S\ref{sec:main}) is a \emph{conditional localisation}: the
existence of such an operator --- on a lift satisfying the
substrate constraints, its positive eigenvalues in
multiplicity-respecting bijection with the upper-half-plane
non-trivial zeros --- implies RH, by the classical argument. The
converse is stated only as the target it is: an open construction
problem. We do not prove RH, and we do not claim an equivalence.

\section{The substrate framework, briefly}
\label{sec:framework}

\subsection{Setup}

\begin{itemize}[leftmargin=*]
  \item $\Vsix$ is the $600$-cell vertex set, $|\Vsix| = 120$,
        identified with the binary icosahedral group $2I$.
  \item $A_1$ is the edge adjacency operator on $\Vsix$.
  \item $\Cphi = (12 + \varphi^{-2}) I - A_1$ is the closure
        operator.
  \item $\Cphi$ has nine eigenvalues in $\Z[\varphi]$ with
        multiplicities summing to $120$.
  \item The Galois action $\sigma: \varphi \mapsto 1 - \varphi$
        partitions the spectrum into $94$ rational (the
        $\tau$-fixed block $V_+$) and $26$ irrational (the
        $\tau$-paired block $V_-$, decomposing further as
        $13 + 13$).
  \item $\Vmin = \ker(\Cphi - \varphi^{-2} I)$ is $1$-dimensional
        and sits inside $V_+$.
\end{itemize}

\subsection{The substrate L-function}

The icosian theta L-function \cite{icosianTriad}:
\[
\Lsub(s) = \zeta_K(s) \cdot \zeta_K(s-1) \cdot C_2(s),
\qquad K = \Q(\sqrt 5),
\]
with $\zeta_K(s) = \zeta(s) \cdot L(s, \chi_5)$ and $C_2(s) = 1 -
2 \cdot 2^{-s} + 2^{2-2s}$.

\subsection{The closure equation}

The triadic closure equation \cite{observerAttractor}:
\[
\Wit = \Phi \times G \times O \quad \text{subject to} \quad C \to 0
\]
has the unique solution $\Wit = \Vmin$, where:
\begin{itemize}[leftmargin=*]
  \item $\Phi$ (Phase) = closure-flow dynamics on $\C^{120}$,
  \item $G$ (Geometry) = $\Vsix$ with $A_1$,
  \item $O$ (Observer) = observer-instance trajectories,
  \item $C = \|P_- v\|^2 \to 0$.
\end{itemize}

% ============================================================
\section{The localisation target}
\label{sec:main}

\subsection{The spectral generator and the witness axis}

\begin{definition}[Substrate lift]
\label{def:lift}
A \emph{substrate lift} of the $\tau$-fixed block is a separable
infinite-dimensional Hilbert space $\mathcal{H}_+$ together with
an isometric embedding $V_+ \hookrightarrow \mathcal{H}_+$ and
extensions of the closure-flow and of the Hecke action of the
Hilbert modular lift of $V_+$ to $\mathcal{H}_+$. No such lift is
constructed in this paper; exhibiting one with the properties
below is the open problem this paper specifies. (The finite block
$V_+$ itself is $94$-dimensional and cannot carry the infinite
spectrum required below --- the lift is not optional.)
\end{definition}

\begin{definition}[Substrate-localised generator]
\label{def:spectral-generator}
A \emph{substrate-localised generator} for the witness invariant
$\Wit = \Vmin$ is a self-adjoint operator $\Twit$ on a
substrate lift $\mathcal{H}_+$ (Definition~\ref{def:lift}) such
that:
\begin{enumerate}[leftmargin=*,label=(\roman*)]
  \item $\Vmin$ is the ground state of $\Twit$ (smallest
        eigenvalue, eigenvalue normalised to $0$ for definiteness);
  \item $\Twit$ commutes with the closure-flow on $V_+$;
  \item $\Twit$ commutes with all Hecke operators acting on the
        Hilbert modular form lift of $V_+$;
  \item the positive spectrum of $\Twit$ is discrete (eigenvalues
        counted with multiplicity; no simplicity is assumed ---
        simplicity of the $\zeta$-zeros is itself open).
\end{enumerate}
\end{definition}

\begin{definition}[Positive spectrum]
\label{def:resonance-depth-set}
For a substrate-localised generator $\Twit$ on $\mathcal{H}_+$,
write
\[
\mathcal{R}(\Twit) = \{t \in \R_+ : \exists v \in
\mathcal{H}_+,\ v \perp \Vmin,\ \Twit v = t \cdot v\}
\]
for its set of strictly positive eigenvalues; whenever this set is
compared with a zero multiset, the comparison is by
multiplicity-respecting bijection (eigenvalue multiplicities
against zero multiplicities), not bare set equality. (In v1 of
this paper this set carried a decorative name; it is simply the
positive
spectrum, and we call it that.)
\end{definition}

\subsection{The conditional statement}

\begin{proposition}[Conditional localisation: the classical direction]
\label{thm:witness-resonance}
Let $\Twit$ be a substrate-localised generator
(Definition~\ref{def:spectral-generator}) on some substrate lift
$\mathcal{H}_+$, and suppose there is a bijection, respecting
multiplicities, between the non-trivial zeros $\rho$ of $\zeta$
with $\mathrm{Im}(\rho) > 0$ and the positive eigenvalues
$t_\rho$ of $\Twit$, such that $\rho = \tfrac12 + i\,t_\rho$
for every such zero. Then RH holds. (The zeros with
$\mathrm{Im}(\rho) < 0$ are determined by the conjugation
symmetry $\rho \mapsto \bar\rho$, so the upper-half-plane set
suffices.)
\end{proposition}

\begin{proof}
Self-adjointness of $\Twit$ forces every eigenvalue $t_\rho$ to
be real, so every upper-half-plane zero satisfies
$\rho = \tfrac12 + i\,t_\rho$ with $t_\rho \in \R$, i.e.\
$\mathrm{Re}(\rho) = \tfrac12$; the conjugation symmetry extends
this to all non-trivial zeros. This is the classical
Hilbert--P\'olya implication; nothing about the substrate is used
in the deduction, and we claim no novelty for it.
\end{proof}

\begin{target}[The substrate localisation target]
\label{target:main}
Construct a substrate lift $\mathcal{H}_+$ and a
substrate-localised generator $\Twit$ on it with
$\mathcal{R}(\Twit) = \{\gamma_n\}$ as in
Proposition~\ref{thm:witness-resonance}. The content of the target
is entirely in the substrate constraints (ground state $\Vmin$,
closure-flow and Hecke commutation): without them the statement
``RH implies such an operator exists'' is trivially true (under RH,
take the diagonal operator on $\ell^2$ whose eigenvalues are the
upper-half-plane ordinates, repeated with multiplicity), and
with them it is open. We do not assert that RH guarantees a
generator satisfying the constraints; that assertion was the
central defect of v1 of this paper and is withdrawn.
\end{target}

\begin{remark}[On the withdrawn equivalence]
v1 stated an equivalence between RH and ``substrate spectral
completeness''. One direction is
Proposition~\ref{thm:witness-resonance}; the other was supported
only by the assumption that the construction can be carried out
under RH, which assumes what it must prove. A third clause
(``the closure equation is spectrally complete'') was definitional
--- equivalent to the first by unwinding the definitions --- and is
retained only as terminology, not as a result. An analogous
conditional statement can be formulated for $L(s,\chi_5)$ on the
$\tau$-paired block $V_-$; it has the same one-direction status
and the same open converse, and we state it below only as a
conjectural target.
\end{remark}

\subsection{What the localisation says precisely}

\begin{equivbox}
The statement of this paper, in one breath:

\medskip
\noindent\itshape
If the witness invariant $\Wit = \Vmin$ admits a
substrate-localised generator on some lift of the $\tau$-fixed
block whose positive eigenvalues are in multiplicity-respecting
bijection with $\zeta$'s upper-half-plane non-trivial zeros (via
$\rho = \tfrac12 + i\,t_\rho$; the conjugate lower half follows
by symmetry), then RH holds. Whether RH conversely
guarantees a generator satisfying the substrate constraints is
open, and is the localisation target.

\medskip
\noindent\upshape
This is not a proof of RH, and not an equivalence. Closing the
target requires either:

\begin{itemize}[leftmargin=*]
  \item \textbf{Bottom-up}: explicitly constructing $\Twit$ from
        the substrate's intrinsic data (closure-flow + cuspidal
        Hilbert modular forms over $\Q(\sqrt 5)$ + $2I$-
        equivariance). This is the Hilbert--P\'olya programme
        applied to V$_{600}$.
  \item \textbf{Top-down}: proving that any
        Hilbert--P\'olya operator on the relevant Hilbert modular
        form space must localise onto $\Vmin$'s span (the
        $\tau$-fixed block); then $\Twit$ exists by analytic
        construction, and RH follows.
\end{itemize}

The first resolution would prove RH
(Proposition~\ref{thm:witness-resonance}); the second would close
the localisation target but still leaves the construction to do.
The substrate framework sharpens where the operator must live; it
does not move the answer.
\end{equivbox}

\subsection{Extension: $L(s, \chi_5)$ on the $\tau$-paired block}
\label{sec:chi5-extension}

\begin{target}[$\tau$-paired analogue, conjectural]
\label{prop:tau-paired-generator}
The same one-direction statement can be formulated for
$L(s,\chi_5)$ on a lift of the $\tau$-paired block $V_-$: a
substrate-localised generator there whose positive eigenvalues are
in multiplicity-respecting bijection with the non-trivial zeros of
$L(s,\chi_5)$ in the upper half-plane (via
$\rho = \tfrac12 + i\,t_\rho$; the conjugate-symmetric lower
half follows) would imply $\mathrm{GRH}_{\chi_5}$ by the
identical self-adjointness argument. The Dedekind
factorisation $\zeta_K = \zeta \cdot L(s,\chi_5)$ motivates the
block assignment. No converse is asserted, and no construction is
given; this is a target, stated for symmetry.
\end{target}

\subsection{Corollaries}

\begin{remark}[Combined conditional statement]
\label{cor:combined}
If substrate-localised generators exist on lifts of both blocks
with the stated spectra, then RH and $\mathrm{GRH}_{\chi_5}$ both
hold, by applying Proposition~\ref{thm:witness-resonance} and
Target~\ref{prop:tau-paired-generator}'s implication to each
factor of $\zeta_K = \zeta \cdot L(s,\chi_5)$. (The finite
direct sum $V_+ \oplus V_- = \C^{120}$ of v1 is replaced by the
direct sum of the lifts; the finite space cannot carry either
spectrum.)
\end{remark}

\section{What this reframe accomplishes and what remains open}
\label{sec:scope}

\subsection{What the reframe accomplishes}

\begin{enumerate}[leftmargin=*]
  \item \textbf{A constraint list for Hilbert--P\'olya, anchored
        in the substrate.} Any operator meeting
        Definition~\ref{def:spectral-generator} --- ground state at
        $\Vmin$, commutation with closure-flow and with the Hecke
        action of the Hilbert modular lift --- would witness RH by
        Proposition~\ref{thm:witness-resonance}. The constraints
        are checkable; the construction is not supplied.
  \item \textbf{An honest size statement.} The $\tau$-fixed
        block itself is $94$-dimensional, so the search is
        \emph{not} finite-dimensional (v1 claimed it was; that
        claim is withdrawn): any candidate lives on an
        infinite-dimensional lift anchored to the block.
  \item \textbf{A precise interface to the closure equation.}
        The closure equation $\Wit = \Phi \times G \times O$
        subject to $C \to 0$ acquires a definite operator-side
        reading --- as a \emph{definition} of what its spectral
        completeness would mean, not as a validated theorem.
\end{enumerate}

\subsection{What remains open}

\begin{enumerate}[leftmargin=*]
  \item \textbf{Constructing $\Twit$ from substrate-intrinsic data.}
        The translation engine has shown
        \cite{translationEngine} that finite-group operators on
        V$_{600}$ alone do not produce $\{\gamma_n\}$; the missing
        ingredient is cuspidal Hilbert modular form data over
        $\Q(\sqrt 5)$. Once such data is available (via LMFDB,
        Magma, or future SAGE), the construction can be attempted.
  \item \textbf{The classical Hilbert--P\'olya conjecture.}
        Independent of the substrate's localisation, the existence
        of a self-adjoint Hilbert--P\'olya operator for $\zeta$ is
        open. The substrate framework restricts the search domain
        but does not produce the operator.
  \item \textbf{Substrate-level positivity arguments.}
        A direct substrate-side proof would need to show that any
        substrate-localised generator on a lift of $V_+$
        \emph{must} have its spectrum on $\R_+$. The current substrate framework permits this
        as a possibility but does not prove it.
\end{enumerate}

\begin{scopebox}
\textbf{What is established.}
\begin{itemize}[leftmargin=*]
  \item One direction (Proposition~\ref{thm:witness-resonance}):
        a substrate-localised generator on a lift, its positive
        eigenvalues in multiplicity-respecting bijection with the
        upper-half-plane non-trivial zeros, would imply
        RH --- by the classical self-adjointness argument, with no
        novelty claimed for the deduction.
  \item A precise constraint list
        (Definition~\ref{def:spectral-generator}) localising where
        such an operator must live.
\end{itemize}

\textbf{What is not established.}
\begin{itemize}[leftmargin=*]
  \item The Riemann Hypothesis, $\mathrm{GRH}_{\chi_5}$, or any
        analytic statement about $\zeta$ or $L(s,\chi_5)$.
  \item The converse: that RH guarantees a generator satisfying
        the substrate constraints (v1's claimed equivalence,
        withdrawn).
  \item The existence of a substrate lift, or of $\Twit$ on any
        lift.
  \item Any new computation: this bundle contains no simulations;
        all structural counts ($120$, $94{+}13{+}13$,
        $\dim\Vmin=1$) are imported from the cited sibling
        artifacts (\texttt{critical-line-pullback},
        \texttt{the-24-600-spectral-bridge}).
\end{itemize}
\end{scopebox}

\section{Conclusion}
\label{sec:conclusion}

The Riemann Hypothesis is localised, not resolved, and not
reformulated as an equivalence. The substrate framework specifies
\emph{where} a Hilbert--P\'olya operator for $\zeta$ would have
to live ($\Vmin$-anchored, on an infinite-dimensional lift of the
$\tau$-fixed block) and \emph{what} structural properties it
would have to satisfy
(Definition~\ref{def:spectral-generator}); its existence with the
stated spectrum would imply RH
(Proposition~\ref{thm:witness-resonance}), by the classical
argument. Whether RH conversely guarantees an operator meeting the
substrate constraints is open --- that is the localisation target,
and the honest name for what this paper contributes is a
\emph{specification}. The analytical step that would close it is
the same one that has been open since Riemann.

\bigskip

\hrule
\bigskip

\noindent
\textbf{Acknowledgements.} This paper synthesises the programme's
preceding work; no new techniques are introduced. The structural
diagram licensed by the closure equation $\Wit = \Phi \times G
\times O$ subject to $C \to 0$ is due to the
\textit{VFD Witness Structure} framework
\cite{observerAttractor}.

\bibliographystyle{plain}
\bibliography{references}

\end{document}
